\chapter{Policy expression reference}
\label{app:policy}

The complete owner policy as deployed. Worker and entry node differ only where
the schedd role requires it; those differences are marked.

\section{Derived quantities}

\begin{lstlisting}
NonCondorLoad  = (TotalLoadAvg - TotalCondorLoadAvg)

LOAD_ADMIT     = (TotalCpus * 0.50)      # entry node: DetectedCpus
LOAD_SUSPEND   = (TotalCpus * 0.75)      # entry node: DetectedCpus

MEM_HARD_FLOOR = ifThenElse((HostMemTotalMB * 0.10) > 1024,
                            (HostMemTotalMB * 0.10), 1024)
MEM_SOFT_FLOOR = (HostMemTotalMB * 0.25)

# Disk floors are fractional on workers and ABSOLUTE on the entry node,
# whose SPOOL shares a mostly-full volume with the OS.
DISK_SOFT_FLOOR = (HostDiskTotalMB * 0.10)     # entry node: 15000
DISK_HARD_FLOOR = (HostDiskTotalMB * 0.05)     # entry node:  8000
\end{lstlisting}

\section{Fit tests}

\begin{lstlisting}
FITS_MEM  = (TARGET.RequestMemory =?= UNDEFINED)
         || (TARGET.RequestMemory < (HostMemAvailMB - MEM_SOFT_FLOOR))
\end{lstlisting}

The \kn{=?=} guard is not optional. A job that specifies no
\kn{request\_memory} yields \kn{UNDEFINED}, which without the guard would make
the whole \kn{\&\&} chain non-true and silently refuse every unspecified job.

\section{The policy}

\begin{lstlisting}
EVICT_PRESSURE = (HostMemAvailMB > 0)
              && (  (HostMemAvailMB  < MEM_HARD_FLOOR)
                 || (HostMemPsiAvg10 > 10.0)
                 || (HostDiskAvailMB < DISK_HARD_FLOOR) )

WANT_SUSPEND = !($(EVICT_PRESSURE))

START    = (NonCondorLoad   < LOAD_ADMIT)
        && (TotalLoadAvg    < TotalCpus * 0.90)
        && (HostCpuPsiAvg10 < 20.0)
        && (HostMemAvailMB  > MEM_SOFT_FLOOR)
        && (HostDiskAvailMB > DISK_SOFT_FLOOR)
        && FITS_CPU && FITS_MEM && FITS_DISK

SUSPEND  = (NonCondorLoad > LOAD_SUSPEND)
CONTINUE = (NonCondorLoad < LOAD_ADMIT)
PREEMPT  = $(EVICT_PRESSURE)

MaxJobRetirementTime = 0
MachineMaxVacateTime = 30
KILL                 = FALSE
\end{lstlisting}

\section{Supporting settings}

\begin{lstlisting}
UPDATE_INTERVAL            = 30      # default 300 - see chapter 3
SCHEDD_INTERVAL            = 30
CGROUP_MEMORY_LIMIT_POLICY = none    # no per-job cap - see chapter 4
JOB_RENICE_INCREMENT       = 10
PREEN_INTERVAL             = 3600    # default 86400 - see chapter 3
PREEN_ARGS                 = -r
\end{lstlisting}

\section{Attributes published by STARTD\_CRON}

Condor publishes none of these itself. Without the cron stanza every expression
above evaluates to \kn{UNDEFINED} and the entire policy is inert
(\S\ref{sec:inert}).

\begin{table}[H]
\centering
\begin{tabularx}{\textwidth}{L{0.30\textwidth} Y}
\toprule
\textbf{Attribute} & \textbf{Source} \\
\midrule
\kn{OwnerSessionActive} & Remote-desktop connections plus active non-greeter graphical sessions \\
\kn{OwnerSessionReason} & Diagnostic string explaining the above \\
\kn{HostMemAvailMB}, \kn{HostMemTotalMB} & \file{/proc/meminfo} \kn{MemAvailable}, not \kn{MemFree} \\
\kn{HostMemPsiAvg10} & \file{/proc/pressure/memory} \\
\kn{HostCpuPsiAvg10} & \file{/proc/pressure/cpu} \\
\kn{HostDiskAvailMB}, \kn{HostDiskTotalMB} & The \kn{EXECUTE} filesystem \\
\kn{HostSpoolAvailMB}, \kn{HostSpoolTotalMB} & The \kn{SPOOL} filesystem \\
\bottomrule
\end{tabularx}
\caption{Attributes the policy depends on, and where each comes from.}
\end{table}
